\documentclass[11pt]{article}
\usepackage[margin=1in]{geometry}
\usepackage{amsmath,amssymb,amsthm,hyperref}
\newtheorem{theorem}{Theorem}
\title{Exact concentration moves, target bridges, and path-dependent barriers}
\author{Joshua Mark}
\date{20 November 2026; research draft, priority unestablished}
\begin{document}\maketitle
Let $\mathcal P_d$ be the probability laws on the sign cube with
zero first and second moments. Write $A$ for the matrix evaluating
the constant, degree-one and degree-two characters, and $b=(1,0,\ldots,0)$.
Thus $\mathcal P_d=\{p\ge0:Ap=b\}$ and $C(p)=\|p\|_2^2$.

\begin{theorem}[An exact monotone support-reduction algorithm]
Given any rational feasible law $p$ with support $S$, repeat the
following operation while the columns $A_S$ are dependent.
Choose nonzero $z\in\ker A_S$, reverse its sign if necessary so
$p\cdot z\ge0$, and set
\[
 t=\min_{j:z_j<0}\frac{p_j}{-z_j},\qquad p\leftarrow p+tz.
\]
Every intermediate law remains nonnegative and exactly feasible,
collision increases strictly at each nontrivial step, and the entire
interpolating path has nondecreasing collision. After at most
$|S|-\operatorname{rank}A_S$ steps it reaches a vertex supported on
at most $1+d+\binom d2$ atoms. All calculations can be rational.
\end{theorem}
\begin{proof}
The constant row of $Az=0$ gives $\sum z_j=0$, so $z$ has both
signs and $t>0$ is finite. The definition of $t$ preserves positivity
and kills at least one atom; moment feasibility follows from $Az=0$.
For $0\le u\le t$,
$C(p+uz)=C(p)+2u(p\cdot z)+u^2\|z\|^2$ is nondecreasing and
strictly increases at the endpoint. The feasible face containing
$p$ in its relative interior has dimension
$|S|-\operatorname{rank}A_S$. Each boundary step passes to a proper
face, decreasing that dimension at least one. At a zero-dimensional
face the active columns are independent. This proves termination
and the support bound. One can choose a minimally supported kernel
vector, making every step a circuit move. This is the elementary
convex-objective endpoint principle familiar from pipage-type
rounding~\cite{AS}, not a claimed new general optimization method.
\end{proof}
The algorithm can terminate at a suboptimal vertex, including the
weighted four-bit examples. It does not guarantee minimal entropy
or a Hadamard construction. The result identifies exactly what local
support deletion can do without a collision barrier.

\begin{theorem}[Sharp barrier for a specified straight bridge]
Let $p,q\in\mathcal P_d$ be distinct and set
$a=C(p)$, $b=C(q)\ge a$, $c=p\cdot q$, $e=a+b-2c>0$.
The feasible bridge $p_t=(1-t)p+tq$ has minimum collision
\[
 \min_{0\le t\le1}C(p_t)
 =a-\frac{(a-c)_+^2}{a+b-2c}.
\]
The minimizer is $t=(a-c)_+/e\le1/2$. Its collision loss is at most
$a^2/(a+b)\le a/2$, with the first bound attained for disjoint
supports. If $c\ge a$ the bridge is monotone and has no loss.
\end{theorem}
\begin{proof}
Expand $C(p_t)=a+2(c-a)t+et^2$. If $c<a$, the minimizer is
$(a-c)/e$, at most one half since $b\ge a$. Since $c\ge0$ and
$c<a\le b$, the function $(a-c)^2/(a+b-2c)$ decreases with $c$:
its derivative is $-2(a-c)(b-c)/(a+b-2c)^2\le0$.
Its value at zero is $a^2/(a+b)$. If $c\ge a$ the quadratic is
nondecreasing on the entire interval.
\end{proof}
If an explicit optimal law $q$ is available, this gives an actual
nonlocal feasible route and a quantitative loss budget. It does not
provide the missing optimal target. Translating $q$ by cube signs
preserves feasibility and collision, and averaging those translates
gives the uniform cube. Consequently some translate has overlap
$p\cdot q\ge2^{-d}$, which may improve the displayed bridge bound.

\section*{A sharp bridge need not be a global barrier}
For $d=5$, use the eight-word linear Walsh code
\[
 C=\{((-1)^{u\cdot m})_{m=1}^5:u\in\mathbb F_2^3\}.
\]
The masks $1,\ldots,5$ include a basis and are distinct and nonzero,
so the code is injective and every coset's uniform law is unbiased
and pairwise independent. There are four disjoint cosets in the
32-word cube. Let $q_1,q_2,q_3$ be three of these uniform laws and
$p=(q_1+q_2)/2$. Then
\[
 C(p)=1/16,\quad C(q_3)=1/8,\quad p\cdot q_3=0.
\]
The straight bridge to $q_3$ has minimum $1/24$ at $t=1/3$ and
collision barrier $1/48$, attaining the sharp general bound.
Nevertheless the bridge from $p$ to $q_1$ has overlap $p\cdot q_1=1/16$
and collision $1/16+t^2/16$, which is monotone. Both targets have
optimal collision $1/8$ by the residue-sensitive bound.
Thus even an exactly sharp barrier to a selected optimum cannot be
promoted to a barrier for reaching \emph{any} optimum. The globally
optimal four-bit barrier required its additional section argument;
there is no automatic extension to larger moment polytopes.

\begin{thebibliography}{9}
\bibitem{AS} A. A. Ageev and M. I. Sviridenko, Pipage rounding:
a new method of constructing algorithms with proven performance
guarantee, Journal of Combinatorial Optimization 8 (2004), 307--328.
\url{https://doi.org/10.1023/B:JOCO.0000038913.96607.c2}.
\end{thebibliography}
\end{document}
